\scanpage{11}
\textcircled{II}\quad\uline{Épinglages et repères}.

Un \uline{repère} d'un polygône combinatoire est
un couple $(s_0,d_0)$ d'un sommet et d'un côté incidents\,; il revient au $=$ de
se donner un couple $(s_0,s_1)$ de deux
sommets adjacents \deleted{\unsure{(ou encore|ou ce qui}} (prendre
$d_0=\{s_0,s_1\}$) ou d'un $s_0\in S$ et de $u\in\omega$
-- sommet et orientation -- (prendre $s_1=us_0$).

\uline{Prop} \uline{1} $\Gamma=\operatorname{Aut}(S,\omega)$ est \inserted{\unsure{simplemt|simplement}} transitif sur l'ens des repères -- \deleted{\unsure{qui est donc}} \inserted{i.e. celui-ci est} un torseur sous $\Gamma$.

C'est en. équivalent à ceci

\uline{Cor} 1 Soient $\Pi=(S,\omega)$, $\Pi'=(S',\omega')$ deux \deleted{\unsure{isom|isom.}} pol
comb \inserted{isomorphes}, et soit $r_0=(s_0,u)$ un repère de
$\Pi$. Alors l'application
\[
\begin{aligned}
\operatorname{Isom}(\Pi,\Pi')&\longrightarrow\operatorname{Rep}(\Pi')\\
\unsure{\theta|\Theta}&\longmapsto\theta(\unsure{r|r_0})
\end{aligned}
\]
est bijective. \unsure{$\vert$}

Un \uline{épinglage} d'un poly. comb. est
la donnée d'un repère (donc \uline{c'est} un
repère) -- \deleted{celui} \inserted{un épinglage (ou repère)} d'un polygône géom.
est un ép. du pol. comb sous-jacent.
Il revient au $=$ de \deleted{dire} se donner un $\Pi_0$-
% archive sheet number: 10

